\documentclass[10pt,a4paper]{amsart}
\usepackage[margin=19mm]{geometry}
\usepackage{amsmath,amssymb,mathtools}
\usepackage[scr=rsfs,cal=euler]{mathalfa}
\usepackage{xfrac}
\usepackage[unicode,hidelinks,bookmarks=false]{hyperref}
\newcommand{\ie}{i.e.\ }
\newcommand{\cf}{cf.\ }
\newcommand{\resp}{respectively\ }
\newcommand{\Subsection}{\subsection}
\providecommand{\nobreakdash}{\nobreak\hbox{-}}
\setlength{\emergencystretch}{2em}
\multlinegap0pt
\usepackage{mdframed}
\usepackage{mdframed}
\usepackage{mdframed}
\usepackage{mdframed}
\usepackage{amssymb}
\usepackage{xcolor}
\usepackage{enumerate}
\usepackage{mathtools}
\usepackage{tikz}
\usepackage{mdframed}
\newtheorem{lemma}{Lemma}
\newcommand{\cL}{\mathcal{L}}
\newcommand{\cM}{\mathcal{M}}
\newcommand{\cT}{\mathcal{T}}
\newcommand{\colorM}{\color{black}}
\newcommand{\gradM}{\operatorname{grad}_{\cM}}
\renewcommand{\v}{{\colorM v}}
\title{Natural gradient descent with momentum}
\author{Anthony Nouy, Agust\'in Somacal}
\date{Source: arXiv:2604.15554v1 (CC BY 4.0)}
\begin{document}
\maketitle

\noindent\emph{Source and licence.} Natural gradient descent with momentum. Version arXiv:2604.15554v1 (CC BY 4.0), distributed by arXiv under the Creative Commons Attribution 4.0 International licence, http://creativecommons.org/licenses/by/4.0/, as recorded in the arXiv licence metadata for this exact version and shown on the abstract page https://arxiv.org/abs/2604.15554v1. Full licence notice: \texttt{licenses/arxiv-2604.15554-CC-BY-4.0.txt}.\par\smallskip

\noindent\textit{Editorial scope.} Counted unit: the article's lemma lem:projection (source0.tex lines 924--932). The article's complete original proof of it is delivered with it (source0.tex lines 933--936, one \texttt{proof} environment of 4 lines). No mathematics is written, repaired or re-derived: the statement and the proof are the article's own lines, in the article's own notation, and no retained line is substituted or re-flowed.  No further source-local context is needed: every cross-reference of every retained line resolves inside the two delivered spans.  Nothing was omitted from any retained span, so the package carries no omission record.  No retained line contains a citation, so the wrapper carries no bibliography and no bibliography entry.  No retained line contains a figure environment, an image-inclusion command or drawing code, so no image is delivered and the package declares no asset dependency.  Editorial formatting only, in addition: an AMS \texttt{amsart} preamble that restates the article's own declarations and macros used by the retained lines, namely the environments \texttt{lemma} on separate counters, together with the standard packages those lines need.  The retained environments print in this document as Lemma 1 in order of appearance, whereas the article prints the same items with the numbering of its own full document.  The complete native source of the article as distributed in the arXiv e-print for this exact version is bundled unchanged as \texttt{sources/198619/source0.tex}.
\begin{lemma}
\label{lem:projection}
Assume $W\subset V$. Then for all $v\in \cM$, $\cL'(v)$ is  a linear form on $W$ and the restriction $\cL'(v)_{\vert W}$ can be identified by Riesz representation with $\nabla _W \cL(v) \in W$. Thus
$$
\gradM \cL(v) = P^W_{\cT_{v}} \nabla_W \cL(\v),
$$
with $ P^W_{\cT_{v}} $ the $W$-orthogonal projection onto $\cT_v.$ 
In particular, for $V$ a Hilbert space and $W=V$, $$\gradM \cL(v) = P^V_{\cT_{v}} \nabla \cL(\v).$$
\end{lemma}

\begin{proof}
By definition of  $\gradM \cL(v)$ and  $\nabla_W \cL(v)$, it holds
$(\gradM \cL(v) , w)_W = \cL'(v)(w) = (\nabla_W \cL(v) , w)_W$ for all $w \in \cT_v$, which  yields $\gradM \cL(v) = P^W_{\cT_v} \nabla_W \cL(\v)$.
\end{proof}

\end{document}
